%%%% CAPÍTULO 1 - INTRODUÇÃO

\chapter{Introdução}\label{cap:introducao}

O futebol inglês possui uma das tradições mais ricas e antigas do cenário
desportivo mundial. A \textit{Premier League}, na sua configuração moderna de
liga independente, foi fundada em 1992 após uma ruptura estratégica com a antiga
\textit{Football League First Division}, cuja fundação original remonta a 1888.
Essa transição visou não apenas uma modernização comercial e técnica, mas também
a consolidação do campeonato como o produto de entretenimento esportivo mais
consumido no planeta~\cite{PremierLeague2024}.

Dentro do contexto competitivo do futebol, o fenômeno conhecido como Fator Casa
(\textit{home advantage}) surge como um dos pilares centrais da cultura do
esporte. Tradicionalmente, assume-se que jogar em domínios próprios confere uma
vantagem competitiva real às equipes mandantes, impulsionada por fatores como o
apoio da torcida, a familiaridade com o gramado, o desgaste logístico imposto ao
visitante e a pressão psicológica exercida sobre a equipe de
arbitragem~\cite{Courneya1992,Nevill1999}.

\section{Problema de Pesquisa}\label{sec:problema}

Embora o Fator Casa seja um conceito amplamente aceito historicamente, o futebol
moderno passou por profundas transformações nas últimas décadas. A globalização
do esporte, a massificação de esquemas táticos altamente estruturados, a melhoria
na logística de viagens e, mais recentemente, a introdução de tecnologias de
auxílio à arbitragem levantaram novos questionamentos.

A pandemia de COVID-19 acrescentou a esse cenário um elemento raro na pesquisa
esportiva: a realização de centenas de partidas com portões fechados, condição
que a literatura passou a tratar como um experimento natural sobre o papel do
público~\cite{Fischer2021}. A interpretação desse episódio, contudo, exige
cautela, pois a ausência de torcida não foi a única alteração do período.

Diante desse cenário, formula-se o seguinte problema de pesquisa: \textbf{até que
ponto o mando de campo ainda se configura como um fator estatisticamente
determinante no desempenho das equipes da \textit{Premier League} na era moderna,
e quais mecanismos sustentam essa vantagem?}

\section{Justificativa}\label{sec:justificativa}

A realização deste estudo justifica-se pela necessidade de validar, por meio de
uma abordagem estritamente quantitativa e empírica, mitos e convenções enraizados
na cultura do futebol profissional. Para além do interesse puramente esportivo, a
análise científica do Fator Casa possui relevância prática direta para analistas
de desempenho, comissões técnicas e para o mercado de ciência de dados aplicada
ao esporte.

Soma-se a isso uma justificativa de ordem metodológica. Análises descritivas do
mando de campo, baseadas em médias de pontos conquistados dentro e fora de casa,
são incapazes de separar a vantagem atribuível ao mando da qualidade dos
adversários enfrentados em cada condição. Modelos estatísticos que estimam
simultaneamente a força de cada clube e um parâmetro de mando permitem essa
separação, e sua aplicação a um recorte recente da \textit{Premier League}
constitui uma contribuição deste trabalho.

\section{Objetivos}\label{sec:objetivos}

Nesta seção são delimitados o objetivo geral e os objetivos específicos que
guiam a execução desta pesquisa.

\subsection{Objetivo Geral}\label{sec:objetivo-geral}

Analisar de forma quantitativa e estatística a influência do mando de campo nos
resultados e nas métricas de desempenho das equipes na \textit{Premier League},
cobrindo o recorte temporal das temporadas de 2015/16 até 2025/26, e estimar essa
influência por meio de modelos que controlem pela força dos adversários.

\subsection{Objetivos Específicos}\label{sec:objetivos-especificos}

Para alcançar o objetivo geral, o trabalho propõe-se a:

\begin{enumerate}
  \item \textbf{Magnitude do domínio:} quantificar a vantagem do mandante em
    pontos, gols e ações ofensivas, reportando estimativas acompanhadas de
    intervalos de confiança e tamanhos de efeito;
  \item \textbf{Heterogeneidade entre equipes:} verificar se equipes que
    terminam a temporada em diferentes faixas da tabela de classificação
    apresentam o mesmo comportamento em relação ao mando de campo;
  \item \textbf{Evolução temporal:} identificar se o Fator Casa apresenta
    tendência de queda, subida ou estabilidade ao longo do período estudado;
  \item \textbf{Papel do público:} avaliar o impacto da ausência de torcida
    sobre a vantagem do mandante, distinguindo o efeito do público de outras
    alterações ocorridas no mesmo período;
  \item \textbf{Modelagem:} estimar o parâmetro de mando de campo por meio de um
    modelo de Poisson bivariado e avaliar a capacidade preditiva de modelos
    estatísticos e de aprendizado de máquina sob validação temporal.
\end{enumerate}

\section{Organização do Trabalho}\label{sec:organizacao}

Este trabalho está estruturado em cinco capítulos. Este primeiro capítulo
apresentou a introdução, o problema de pesquisa, a justificativa e os objetivos
do estudo. O Capítulo~\ref{cap:fundamentacao} aborda a fundamentação teórica,
trazendo o levantamento da literatura sobre o Fator Casa, os conceitos da
metodologia CRISP-DM, os modelos estatísticos aplicados a resultados de futebol e
os critérios de avaliação de modelos probabilísticos. O
Capítulo~\ref{cap:metodologia} detalha os materiais e métodos, descrevendo a
coleta de dados, as tecnologias utilizadas, as técnicas de engenharia de dados
aplicadas e os procedimentos adotados para garantir a reprodutibilidade da
análise. O Capítulo~\ref{cap:resultados} apresenta e discute os resultados
obtidos. Por fim, o Capítulo~\ref{cap:conclusao} apresenta as conclusões, as
limitações do estudo e as possibilidades de trabalhos futuros.
